\section{Does promoting the hidden torsion make marginal Fisher--Rao useful?}
\label{sec:cv_does_it_help}

We can now compare the three coordinates on a common axis: the one-dimensional torsion
$\varphi_1$ (\cref{sec:case_pentane}), the harder one-dimensional distance $R_{15}$
(\cref{sec:cv_harder_1d}), and the two-dimensional torsion $(\varphi_1,\varphi_2)$
(\cref{sec:cv_2d}). The organizing question is whether mFR's benefit tracks \emph{measured
starvation of the biased coordinate}, and whether any apparent gain is a genuine
support gain or a lineage-diversity loss.

\subsection{Method comparison}
On a cell the ABF screen flags as \emph{starved} the pre-registered gate calls for the full
ladder --- multi-replica shared-bias ABF, ABF+mFR with estimated / uniform / oracle targets,
and an independently swept OPES --- which is what \cref{sec:cv_harder_1d} reports for
$R_{15}$. Since no two-dimensional cell is starved, the gate instead calls for the compact
\emph{neutrality control} reported here: ABF, the deployable estimated-target mFR, a labelled
aggressive probe, and OPES (the uniform/oracle diagnostic targets are informative only where
mFR has a support deficit to repair, and on $R_{15}$ they proved indistinguishable from the
estimated target). Both are run at matched seeds, initial conditions, replicas, steps and
force-evaluation budget. We are explicit about the asymmetry in how hyperparameters were set,
because it is not the same on the two arms. On $R_{15}$ the mFR rate was chosen from a ladder
run on disjoint held-out tuning seeds, and the OPES baseline was swept independently over
barrier, kernel width and pace. On the two-dimensional arm neither search was run: the mFR
rate ($0.01$) and the OPES configuration (barrier $8$, $\sigma=0.30$, pace $400$) were fixed
\emph{a priori} in the configuration and never tuned, because the pre-registered gate does not
authorise a hyperparameter study on a cell no screen classifies as starved. That is a real
limitation of the two-dimensional OPES number in particular, which should be read as ``OPES at
a plausible default'' rather than ``OPES at its best''. It does not weaken the mFR verdict,
which is a null: an untuned mFR that lands within $0.2\%$ of ABF cannot have been flattered by
the absence of tuning. The principal causal contrast is shared-bias ABF versus
ABF+\emph{estimated-target} mFR. \Cref{tab:cv_2d_methods} collects the
two-dimensional production comparison and \cref{fig_cv_methods} the matched-seed effect sizes
and genealogy.

The deployable 2-D mFR verdict on the primary cell is \textbf{\CvTwoDmFRVerdict}
(median change in final $L_2(F)$ vs ABF: \CvTwoDmFRRelFinal\,\%): ABF reaches $0.724$ and mFR
$0.725$, winning $0/6$ matched seeds, with every genealogy, event-cap and basin guard satisfied
--- a genuine null, not a masked false improvement. A \emph{positive} result requires all
twelve pre-registered criteria; a lower free-energy error accompanied by degraded basin
occupancy, conditional structure or genealogy beyond the stated bounds is a \emph{false}
improvement, not a success. OPES at its default configuration is again far worse ($1.682$,
$+132\%$); as noted above it was not swept here, so this is a weaker statement than the
independently tuned $R_{15}$ comparison.

Comparing the two coordinates exposes a second regularity, about the \emph{cost} rather than
the benefit of birth--death. The aggressive probe (rate $0.10$) costs only $+0.9\%$ here, with
ancestor ESS holding at $0.35\,N$, whereas the same probe on the starved $R_{15}$ cell cost
$+34\%$ with ESS collapsing to $0.06\,N$. The difference is how fast the system re-equilibrates
after a clone: on the torus ABF sustains ${\sim}130$ joint round trips per replica, so cloned
configurations decorrelate almost immediately and aggressive reallocation is merely useless;
on the barrier-limited distance coordinate the clones stay correlated, so the same mechanism
actively destroys the mean-force estimate. Birth--death is therefore cheap where it is
unnecessary and expensive precisely where one would want to use it.

\paragraph{Frozen-bias control: the same test, the opposite sign.} Freezing the learned
two-dimensional bias $B(\varphi_1,\varphi_2)$, running fresh dynamics with no updates and no
birth--death, and reconstructing $F=B-\beta^{-1}\log p_B+C$ reproduces the reference far
\emph{better} than the online mean-force integral: $0.378$ against $0.722$ for ABF and $0.378$
against $0.723$ for mFR, i.e.\ the online residual overstates the error in the learned bias by
about a factor of two. This completes a three-way control with a consistent reading. Where the
screen calls the coordinate easy --- the $\varphi_1$ torsion ($0.05$ vs $0.27$) and the
two-dimensional torus here --- freezing \emph{improves} the reconstruction, so the online
residual is a mean-force estimator-resolution floor and there is no systematic error for mFR
to remove. Where the screen calls the coordinate starved --- $R_{15}$ --- freezing makes it
\emph{worse} ($1.92$ vs $1.51$; \cref{sec:cv_harder_1d}), because the learned bias is itself
built from deficient sampling. The control therefore discriminates the two situations
independently of the starvation screen, and agrees with it in all three cases. That the two
methods' reconstructions here are identical to three decimals ($0.3778$ vs $0.3777$) is a
further statement of the null: mFR's learned bias is neither better nor worse than ABF's.
% auto-generated
\begin{table}[t]\centering
\caption{Two-dimensional production comparison (median over seeds).}\label{tab:cv_2d_methods}
\begin{tabular}{llrrrr}
\toprule
cell & method & $L_2(F)$ & basins & cond.\ TV & anc.\ ESS/N\\
\midrule
2D $\beta{=}2$ TT & abf & 0.724 & 9 & 0.389 & --\\
2D $\beta{=}2$ TT & fr active & 0.730 & 9 & 0.404 & 0.35\\
2D $\beta{=}2$ TT & fr estimated & 0.725 & 9 & 0.390 & 0.84\\
2D $\beta{=}2$ TT & opes & 1.682 & 9 & 0.268 & --\\
\bottomrule
\end{tabular}
\end{table}


\begin{figure}[t]\centering
  \includegraphics[width=0.98\textwidth]{fig_cv_04_methods_comparison.png}
  \caption{Matched-seed relative change in final $L_2(F)$ of each method vs ABF (left; the
  shaded band is the $\pm10\%$ practical-equivalence margin, the dashed line the $-15\%$
  positive threshold) and the final ancestor ESS of each mFR configuration (right; the dashed
  line is the $0.30\,N$ genealogy floor).}
  \label{fig_cv_methods}
\end{figure}

\begin{figure}[t]\centering
  \includegraphics[width=0.9\textwidth]{fig_cv_05_fr_birth_death.png}
  \caption{Spatial birth (clone-source) and death maps of the 2-D mFR birth--death process on
  the torsion torus, showing where reallocation moves probability mass.}
  \label{fig_cv_birthdeath}
\end{figure}

\subsection{Cross-case synthesis}
\Cref{tab:cv_starvation_map} places all biased coordinates on one starvation map. The pattern
across the study is that mFR's usefulness is gated by \emph{marginal starvation of the biased
coordinate}: where ABF already resolves the coordinate (the $\varphi_1$ marginals, and any
coordinate the screen classifies easy), mFR has no support deficit to repair and is at best
practically equivalent; the residual difficulty then lives in an \emph{orthogonal} hidden
coordinate that mFR --- acting only on the biased marginal --- cannot drive. Promoting the
hidden torsion into the CV removes the orthogonal-coordinate excuse, so the two-dimensional
comparison isolates whether mFR reallocation on the \emph{actual} slow manifold helps, and
distinguishes a support gain from a lineage-diversity loss (ancestor ESS) and from
two-dimensional density-estimation noise (empty-cell score spikes, curl before projection).
% auto-generated
\begin{table}[t]\centering
\caption{Cross-case ABF starvation synthesis along the biased coordinate.}\label{tab:cv_starvation_map}
\begin{tabular}{lll}
\toprule
coordinate & molecule & ABF starved?\\
\midrule
$\varphi_1$ (1-D dihedral) & butane/pentane & no (prior work)\\
$R_{15}$ (1-D distance) & pentane & \CvRfifteenStarv\\
$R_{14}$ (1-D distance) & butane & \CvRfourteenStarv\\
$(\varphi_1,\varphi_2)$ (2-D) & pentane & \CvTwoDStarvPrimary\ (primary)\\
\bottomrule
\end{tabular}
\end{table}


The three coordinates answer Q3 cleanly, and the answer is not the same in each. For the
one-dimensional torsion $\varphi_1$ and for the two-dimensional torus CV, the limiting factor
is \textbf{(i) lack of marginal starvation}: ABF resolves the biased coordinate to the
estimator floor (in 2-D it biases both dihedrals directly, discovering all nine basins and
mixing strongly even cold and localized), so mFR has no support deficit to repair --- its
measured verdict on the primary 2-D cell is \CvTwoDmFRVerdict. The density-estimation machinery is \emph{not} the bottleneck --- the Gram matrix is
well conditioned, the FFT Poisson projection removes real curl, and the decoupled gate
recovers $V_4\!\oplus V_4$ to the shared floor --- so we can exclude (ii). For the harder
distance CV $R_{15}$, ABF \emph{is} starved, which isolates the decisive distinction: the
starvation is a \textbf{dynamical bottleneck}, not a support deficit reallocation can fill
with useful samples. Extended replicas rarely cross the cold torsional barriers that change
$R_{15}$, and mFR can only clone those that already have. The support diagnostics show the
resulting trade explicitly: pushed hard enough to act, mFR \emph{does} repair the nominal
support (under-supported bin fraction $0.223\to0.082$) yet makes the free energy $34\%$
\emph{worse}, because the newly occupied bins hold clones rather than independent samples
(ancestor ESS $\to0.06\,N$, and round trips per replica \emph{falling} $8.9\to6.0$). The
limiting factors are therefore \textbf{(iii) poor discovery} --- amplification is not
discovery --- inseparably compounded by \textbf{(v) genealogical collapse}: reallocation here
converts a support deficit into a diversity deficit, improving occupancy while degrading the
effective sample size the mean-force estimator actually consumes. Conditional/orthogonal relaxation (iv) is a smaller effect here, since the
hidden torsion conditional $p(\varphi_1,\varphi_2\mid R)$ stays well captured. The unifying
statement is that marginal starvation of the biased coordinate is \emph{necessary but not
sufficient} for mFR to help: it also requires that the starvation be a support deficit its
whole-configuration reallocation can fill, rather than a barrier-crossing bottleneck it can
only amplify. In every molecular collective variable we tried --- easy torsion, hard distance,
two-dimensional torsion --- at least one of these preconditions fails, and the matched-seed
verdicts reported above follow accordingly. The diagnostic oracle target and the frozen-bias
control are reported alongside to separate the accuracy of the learned bias from the shared
online mean-force resolution floor, so that a residual online $L_2$ common to every method is
not mistaken for a systematic error mFR could have removed.
